%% ============================================================
%% Appendix A: Inner Products, Norms, and Function Spaces
%% ============================================================

\chapter{Inner Products, Norms, and Function Spaces}
\label{app:innerproduct}

This appendix provides the minimum background in functional analysis
needed to read Appendix~\ref{app:adjointformalism} and Smith's companion
volume~\cite{smith2014uncertainty}. Readers who are comfortable with inner products
and the $L^2$ space can proceed directly to Appendix~\ref{app:adjointformalism}.

%%------------------------------------------------------------
\section{Inner Products and Norms}
\label{appsec:innerproducts}
%%------------------------------------------------------------

An \textbf{inner product}\index{inner product} on a vector space $V$ over $\R$ is a function
$\langle \cdot, \cdot \rangle : V \times V \to \R$ satisfying:
\begin{enumerate}[nosep]
  \item \textit{Symmetry}: $\langle u, v \rangle = \langle v, u \rangle$.
  \item \textit{Linearity in the first argument}:
        $\langle \alpha u + \beta w, v \rangle
        = \alpha\langle u, v\rangle + \beta\langle w, v\rangle$.
  \item \textit{Positive definiteness}: $\langle u, u \rangle \geq 0$,
        with equality only when $u = 0$.
\end{enumerate}

The \textbf{induced norm}\index{norm!induced} is $\|u\| = \sqrt{\langle u, u \rangle}$.
The Cauchy-Schwarz inequality states $|\langle u, v \rangle| \leq \|u\|\,\|v\|$.

\medskip

\noindent\textbf{Example~A.1 (Euclidean inner product).}
For $\vec{u}, \vec{v} \in \R^n$:
\[
  \langle \vec{u}, \vec{v} \rangle = \vec{u}^T\vec{v} = \sum_{i=1}^n u_i v_i.
\]
This is the dot product. The induced norm is the Euclidean norm
$\|\vec{u}\|_2 = \sqrt{\sum u_i^2}$.

\noindent\textbf{Example~A.2 (Weighted inner product).}
For a symmetric positive definite matrix $W \in \R^{n\times n}$:
$\langle \vec{u}, \vec{v}\rangle_W = \vec{u}^T W \vec{v}$.
The induced norm $\|\vec{u}\|_W = \sqrt{\vec{u}^T W \vec{u}}$ is the
$W$-weighted Euclidean norm. SA computations sometimes use this when
parameters have very different scales.

%%------------------------------------------------------------
\section{The $L^2$ Inner Product}
\label{appsec:L2}
%%------------------------------------------------------------

For functions $f, g : [a,b] \to \R$ that are square-integrable
(i.e., $\int_a^b f(t)^2\,dt < \infty$), the \textbf{$L^2$ inner product}\index{$L^2$ inner product} is:
\begin{equation}
  \langle f, g \rangle_{L^2} = \int_a^b f(t)\,g(t)\,dt.
  \label{eq:L2ip}
\end{equation}
The induced norm is $\|f\|_{L^2} = \sqrt{\int_a^b f(t)^2\,dt}$.

The space $L^2([a,b])$ of all square-integrable functions on $[a,b]$
with inner product~\eqref{eq:L2ip} is a \textbf{Hilbert space}\index{Hilbert space}:
a complete inner product space.
Completeness means that every Cauchy sequence converges to a function
in $L^2([a,b])$.

For Sobol decomposition (Chapter~9), the relevant space is $L^2([0,1]^m)$
with inner product $\langle f, g\rangle = \int_{[0,1]^m} f(\vec{q})g(\vec{q})\,d\vec{q}$.

%%------------------------------------------------------------
\section{Integration by Parts as Transposition}
\label{appsec:IBP}
%%------------------------------------------------------------

The key connection between adjoint sensitivity analysis and linear algebra
is this: integration by parts is the function-space analog of the
matrix transpose.

In $\R^n$, for vectors $\vec{u}$, $\vec{v}$ and a matrix $A$:
\[
  \langle A\vec{u}, \vec{v}\rangle = (A\vec{u})^T\vec{v} = \vec{u}^T(A^T\vec{v})
  = \langle\vec{u}, A^T\vec{v}\rangle.
\]
So $A^T$ is the operator satisfying $\langle Au, v\rangle = \langle u, A^Tv\rangle$.

In $L^2([0,T])$, for functions $u, v$ and the differentiation operator $L = d/dt$:
\[
  \langle Lu, v\rangle_{L^2}
  = \int_0^T \dot{u}(t)\,v(t)\,dt
  = \bigl[u(t)v(t)\bigr]_0^T - \int_0^T u(t)\,\dot{v}(t)\,dt
  = \bigl[uv\bigr]_0^T + \langle u, L^*v\rangle_{L^2},
\]
where $L^* = -d/dt$ (with appropriate boundary conditions) is the
\textbf{adjoint operator} of $L = d/dt$.
The boundary term $[uv]_0^T$ gives rise to the terminal condition
$\lambda(T) = -h'(u(T))$ in the adjoint ODE derivation of Chapter~6.

This correspondence makes precise what we stated informally throughout
the book: the adjoint ODE is the ODE whose differential operator
is the transpose of the forward ODE's operator, exactly as $A^T$
is the transpose of $A$.

%%------------------------------------------------------------
\section{Orthogonality and the ANOVA Decomposition}
\label{appsec:ortho}
%%------------------------------------------------------------

Two elements $u, v$ of an inner product space are \textbf{orthogonal}\index{orthogonality}
if $\langle u, v\rangle = 0$.

In $L^2([0,1]^m)$ with the uniform measure, the functions appearing in
the Sobol ANOVA decomposition (Chapter~9) are mutually orthogonal.
Specifically, for the decomposition
$f = f_0 + \sum_i f_i + \sum_{i<j} f_{ij} + \cdots$,
each term is orthogonal to every other term:
$\langle f_{u}, f_{v}\rangle_{L^2} = 0$ when $u \neq v$ (as index sets).

This orthogonality is why the output variance decomposes as a sum:
$\Var(Y) = \sum_i D_i + \sum_{i<j} D_{ij} + \cdots$.
It also explains why the decomposition requires independent parameters:
independence of the $Q_i$ makes the integration over $[0,1]^m$ factorize,
which is what produces the orthogonality.

%%------------------------------------------------------------
\section{Exercises}
\label{appsec:A_exercises}
%%------------------------------------------------------------

\begin{selfcheck}
Let $f(t) = e^{-t}$ and $g(t) = \cos(\pi t)$ on $[0, 1]$.
(a) Compute $\langle f, g \rangle_{L^2}$ numerically (use MATLAB's
\texttt{integral} function).
(b) Are $f$ and $g$ orthogonal in $L^2([0,1])$?
(c) Compute $\|f\|_{L^2}$ and $\|g\|_{L^2}$, and verify the
Cauchy--Schwarz inequality $|\langle f,g\rangle| \leq \|f\|\,\|g\|$.
\end{selfcheck}
\begin{scAnswer}
(a) $\langle f, g\rangle_{L^2} = \int_0^1 e^{-t}\cos(\pi t)\,dt$.
Integrating by parts twice: result $\approx 0.218$ (verify numerically).
(b) No — the inner product is nonzero, so $f$ and $g$ are not orthogonal.
(c) $\|f\|_{L^2}^2 = \int_0^1 e^{-2t}\,dt = (1-e^{-2})/2 \approx 0.432$,
so $\|f\|_{L^2} \approx 0.657$.
$\|g\|_{L^2}^2 = \int_0^1 \cos^2(\pi t)\,dt = 1/2$, so $\|g\|_{L^2} = 1/\sqrt{2} \approx 0.707$.
Cauchy--Schwarz: $|0.218| \leq 0.657 \times 0.707 \approx 0.464$. \checkmark
\end{scAnswer}

\begin{selfcheck}
The integration-by-parts identity in Section~\ref{appsec:IBP}
establishes that $L^* = -d/dt$ (with zero terminal condition)
is the adjoint of $L = d/dt$ (with zero initial condition)
in $L^2([0,T])$.
(a) Write out $\langle Lu, v\rangle_{L^2}$ explicitly for
$u(t) = t^2$ and $v(t) = t$ on $[0,1]$.
(b) Compute $\langle u, L^*v\rangle_{L^2}$ for the same functions.
(c) Verify that the two results differ by the boundary term $[u(t)v(t)]_0^1$.
\end{selfcheck}
\begin{scAnswer}
(a) $Lu = 2t$. $\langle Lu, v\rangle = \int_0^1 2t \cdot t\,dt = \int_0^1 2t^2\,dt = 2/3$.
(b) $L^*v = -dv/dt = -1$. $\langle u, L^*v\rangle = \int_0^1 t^2 \cdot(-1)\,dt = -1/3$.
(c) Boundary term: $[u(t)v(t)]_0^1 = 1\cdot1 - 0\cdot0 = 1$.
Check: $\langle Lu,v\rangle = \langle u, L^*v\rangle + [uv]_0^1 \Rightarrow 2/3 = -1/3 + 1 = 2/3$. \checkmark
\end{scAnswer}

\begin{selfcheck}
The ANOVA decomposition in $L^2([0,1]^2)$ for the model
$f(Q_1, Q_2) = Q_1 + Q_2 + Q_1 Q_2$ (with $Q_i \sim \mathcal{U}(0,1)$) is:
$f = f_0 + f_1(Q_1) + f_2(Q_2) + f_{12}(Q_1,Q_2)$.
Compute each term and verify that all terms are mutually orthogonal
in $L^2([0,1]^2)$.
\end{selfcheck}
\begin{scAnswer}
$f_0 = \mathbb{E}[f] = 1/2 + 1/2 + 1/4 = 5/4$.
$f_1(Q_1) = \mathbb{E}[f|Q_1] - f_0 = Q_1 + 1/2 + Q_1/2 - 5/4 = 3Q_1/2 - 3/4$.
$f_2(Q_2) = 3Q_2/2 - 3/4$ (by symmetry).
$f_{12} = f - f_0 - f_1 - f_2 = Q_1Q_2 - Q_1/2 - Q_2/2 + 1/4 = (Q_1-1/2)(Q_2-1/2)$.
Orthogonality: $\langle f_1, f_2\rangle = \int_0^1\int_0^1 (3Q_1/2-3/4)(3Q_2/2-3/4)\,dQ_1\,dQ_2
= \left[\int_0^1(3Q/2-3/4)\,dQ\right]^2 = 0^2 = 0$. \checkmark
Similarly all other pairs integrate to zero.
\end{scAnswer}
